\section{Diseño experimental y reproducibilidad}

Los experimentos del capítulo se generan de forma reproducible a partir del código incorporado en el repositorio. Esta sección resume los parámetros, el modelo de ruido, las métricas y las verificaciones automáticas usadas para producir las figuras y tablas.

\subsection*{Parámetros canónicos}

La configuración numérica es fija en todos los experimentos:
\[
N=256,
\qquad
R=0.9,
\qquad
N_t=256,
\qquad
N_{\mathrm{fft}}=1024,
\]
\[
M=90,
\qquad
\Delta\theta=2^\circ,
\qquad
\texttt{circle=True}.
\]
La opción \texttt{circle=True} impone la hipótesis de soporte en el círculo inscrito y mantiene $N_t=256$ muestras de detector; no modifica la geometría de adquisición de haz paralelo. La longitud $N_{\mathrm{fft}}=1024$ se usa solo como longitud auxiliar para el filtrado por FFT con zero-padding centrado.\\

El código de reproducción se ejecuta desde la raíz del repositorio mediante
\[
\texttt{python code/generate\_chapter3.py}.
\]
Las versiones de NumPy, SciPy, scikit-image y Matplotlib están fijadas en el archivo de requisitos del código. Los detalles de cada corrida se registran en los metadatos de resultados, y la auditoría de normalización queda guardada junto con esos resultados.

\subsection*{Ruido}

Para evaluar estabilidad se añade ruido gaussiano aditivo al sinograma:
\[
g_{\mathrm{ruido}}=g+\eta,
\qquad
\eta_{\ell,k}\sim\mathcal N(0,\sigma_{\mathrm{noise}}^2),
\]
con
\[
\sigma_{\mathrm{noise}}=0.05\max |g|.
\]
La semilla aleatoria es $42$, y la misma realización ruidosa se usa para todos los filtros. Por tanto, las diferencias observadas bajo ruido se deben al filtrado y no a cambios en la muestra aleatoria.

\subsection*{Métricas}

Las métricas se calculan únicamente sobre los índices correspondientes al disco de soporte del fantoma:
\[
\mathcal I_R=\{(i,j):x_i^2+y_j^2\le R^2\}.
\]
Si $u$ es el fantoma y $v$ una reconstrucción, se usa
\[
\mathrm{RMSE}(u,v)
=
\left(
\frac{1}{\#\mathcal I_R}
\sum_{(i,j)\in\mathcal I_R}
(u_{i,j}-v_{i,j})^2
\right)^{1/2},
\]
y el error relativo en norma discreta $\ell^2$,
\[
\mathrm{Err}_{\ell^2}(u,v)
=
\frac{
\left(\sum_{(i,j)\in\mathcal I_R}(u_{i,j}-v_{i,j})^2\right)^{1/2}
}{
\left(\sum_{(i,j)\in\mathcal I_R}u_{i,j}^2\right)^{1/2}
}.
\]
La tabla de resultados también reporta la razón de medias entre reconstrucción y fantoma sobre $\mathcal I_R$. Esta cantidad se conserva solo como diagnóstico de amplitud; no es una métrica principal ni se utiliza como factor de escala.

\subsection*{Verificaciones automáticas}

La generación comprueba dimensiones de imagen, sinograma y multiplicadores; ausencia de valores no finitos; paridad de los símbolos; valores en frecuencia cero; aplicación independiente por columnas; orientación espacial; y reproducibilidad del hash de métricas al ejecutar dos veces con la misma semilla. No se normalizan las reconstrucciones contra el fantoma ni se aplica corrección empírica de amplitud.
